\chapter{BIODATA TIM PENELITI}

\section*{1. Ketua Peneliti}

\textbf{A. Identitas Diri}

{\footnotesize
\begin{longtable}{@{}p{0.8cm}p{6cm}p{0.3cm}p{8cm}}
1. & Nama Lengkap & : & Dian Hanifudin Subhi, S.Kom., M.Kom. \\
2. & Jenis Kelamin & : & Laki-laki \\
3. & Jabatan Fungsional & : & Asisten Ahli \\
4. & NIP/NIK/Identitas Lainnya & : & 198806102019031018 \\
5. & NIDN & : & 0010068807 \\
6. & Tempat dan Tanggal Lahir & : & Gresik, 10 Juni 1988 \\
7. & E-mail & : & dhanifudin@polinema.ac.id \\
8. & Nomor Telepon/HP & : & 082230743546 \\
9. & Alamat Kantor & : & Jl. Soekarno-Hatta 9 Malang \\
10. & Nomor Telepon/Faks & : & 0341-404424/0341-404420 \\
11. & Lulusan yang Telah Dihasilkan & : & D3=12; D4=35; S1= - ; S2= - ; S3= - \\
12. & Mata Kuliah yang diampu & : & Praktikum Dasar Pemrograman, Dasar Pemrograman, Struktur Data, Praktikum Struktur Data, Pemrograman Berbasis Objek, Pemrograman Mobile, Pemrograman Web Lanjut, Pemrograman Web Framework, Komputasi Awan, Big Data, Proyek Sistem Informasi, Penjaminan Mutu Perangkat Lunak, Analisis Desain Berorientasi Obyek \\
\end{longtable}
}

\textbf{B. Riwayat Pendidikan}

{\footnotesize
\begin{longtable}{|p{3.5cm}|p{5.5cm}|p{5.5cm}|}
\hline
& \textbf{S-1} & \textbf{S-2} \\
\hline
Nama Perguruan Tinggi & Institut Teknologi Sepuluh Nopember & Institut Teknologi Sepuluh Nopember \\
\hline
Bidang Ilmu & Teknik Informatika & Teknik Informatika \\
\hline
Tahun Masuk - Lulus & 2006 – 2011 & 2011-2015 \\
\hline
Judul Skripsi/Thesis/Disertasi & Desain dan Implementasi Multiplayer Word Game Generator dan Solver dengan Menggunakan Plurk API & Efisiensi Tracking Multi Target dengan Model Interaksi Publish-Subscribe Adaptif pada Lingkungan Bergerak \\
\hline
Nama Pembimbing & Imam Kuswardayan, S.Kom., MT. & Waskitho Wibisono, S.Kom., M.Eng., Ph.D. \\
\hline
\end{longtable}
}

\textbf{C. Pengalaman Penelitian dalam 5 Tahun Terakhir}

{\footnotesize
\begin{longtable}{|c|c|p{7cm}|c|}
\hline
\textbf{No.} & \textbf{Tahun} & \textbf{Judul Penelitian} & \textbf{Sumber} \\
\hline
1 & 2021 & Penerapan Sql Similarity dan Multi-Instance Mysql Database untuk Proses Penilaian Otomatis pada Sistem Automated Assessment & DIPA \\
\hline
2 & 2022 & Sistem Pemantauan Ketersediaan dan Peretasan Web Terintegrasi dengan Chatbot & DIPA \\
\hline
3 & 2024 & Evaluasi Penilaian Otomatis Pemrograman Web Laravel pada Platform LAIBA (Learning Application in Balanced Assessment) & DIPA \\
\hline
4 & 2025 & Optimalisasi Layanan Aftersales untuk Efisiensi pada Aplikasi Rekachain di PT Rekaindo Global Jasa & DIPA \\
\hline
5 & 2025 & Implementasi Digitalisasi Memo dan Risalah Rapat Berbasis Website dengan Approval Design Pattern untuk Meningkatkan Efisiensi Dan Akurasi Proses Persetujuan di PT. Rekaindo Global Jasa & DIPA \\
\hline
\end{longtable}
}

\textbf{D. Pengalaman Pengabdian Kepada Masyarakat dalam 5 Tahun Terakhir}

{\footnotesize
\begin{longtable}{|c|c|p{7cm}|c|}
\hline
\textbf{No.} & \textbf{Tahun} & \textbf{Judul Pengabdian} & \textbf{Sumber} \\
\hline
1 & 2020 & Rancang Bangun Web Profile di LPK Pusat Terapi A Plus & DIPA \\
\hline
2 & 2020 & Pelatihan Penerapan Aplikasi Point of Sales pada Rumah Herbal Ahsana Malang & DIPA \\
\hline
3 & 2021 & Pengembangan Point of Sales Sebagai Usaha Memaksimalkan Kegiatan Bisnis UMKM & DIPA \\
\hline
4 & 2022 & Pelatihan Kemasan Produk Anyaman Bambu untuk Pengaruh Display Produk dalam Promosi Terhadap Pembelian Konsumen di Desa Duwet Kecamatan Tumpang Kabupaten Malang & DIPA \\
\hline
5 & 2022 & Pelatihan Penggunaan Aplikasi Penyimpanan Arsip Digital Pribadi dan Keluarga & DIPA \\
\hline
6 & 2022 & Pengelolaan Kandang Pembesaran Anak Ayam Petelor desa Sumberjo Kulon Kec Ngunut Kab Tulungagung & DIPA \\
\hline
7 & 2024 & Pelatihan dan Pembuatan Brosur untuk Meningkatkan Penjualan Produk UMKM Desa Wringinsongo Kecamatan Tumpang Kabupaten Malang & DIPA \\
\hline
8 & 2025 & Penerapan Ipteks Dalam Pengembangan Sistem Informasi Untuk Meningkatkan Promosi Produk Unggulan Paguyuban Petani Mawar Desa Gunungsari, Kota Batu & DIPA \\
\hline
9 & 2025 & Peningkatan Efisiensi Produksi Melalui Penerapan Sistem Informasi Peramalan Kebutuhan Kacang Kedelai Di Pabrik Tahu Melati, Batu & DIPA \\
\hline
\end{longtable}
}

\textbf{E. Pengalaman Penulisan Artikel Ilmiah dalam Jurnal dalam 5 Tahun Terakhir}

{\footnotesize
\begin{longtable}{|c|p{6cm}|p{3.5cm}|p{4cm}|}
\hline
\textbf{No.} & \textbf{Judul Artikel Ilmiah} & \textbf{Nama Jurnal} & \textbf{Volume/Nomor/Tahun} \\
\hline
1 & Sistem Kontrol dan Pemberitahuan Status Peladen Jarak Jauh Berbasis Chatbot & Jurnal Teknik Ilmu dan Aplikasi 2 (1), 20-23 & Volume 2 No 1 2021 \\
\hline
2 & Pengelolaan Arsip Digital di BUMDES Desa Duwet Kecamatan Tumpang Kabupaten Malang Menggunakan Teknologi Informasi & Jurnal Pengabdian kepada Masyarakat 2 (4) – 4627-4632 & Volume 2 No 4 2022 \\
\hline
3 & Penerapan Metode WASPAS pada Sistem Pendukung Keputusan Pemilihan Perusahaan Sekuritas & Journal of Multidisciplinary Inquiry in Science, Technology and Educational Research 1 (4) 1950-1964 & Volume 1 No 4 \\
\hline
4 & Sistem Penentu Penerimaan Bantuan Sosial (BLT-DD) untuk Keluarga Kurang Mampu & Jurnal Informatika Polinema 11 (1) 75-82 & Volume 11 No 1 \\
\hline
5 & Implementasi Clean Architecture pada Sistem Manajemen Tiket Penanganan Insiden & Jurnal Informatika Polinema 12 (1) 195-202 & Volume 12 No 1/2025 \\
\hline
6 & Evaluasi Penilaian Otomatis Pemrograman Web Laravel pada Platform LAIBA & Jurnal Minfo Polgan 14 (2) 2685-2693 & Volume 14 No 2/2025 \\
\hline
7 & Inovasi Pengelolaan Lomba Pramuka Berbasis Website dalam Rangka Penerapan Konsep Smart School & Media Pengabdian kepada Masyarakat 11 (3), 326 - 340 & 2025 \\
\hline
\end{longtable}
}

\textbf{F. Karya Buku dalam 5 Tahun Terakhir}

{\footnotesize
\begin{longtable}{|c|p{5cm}|c|c|p{3cm}|}
\hline
\textbf{No.} & \textbf{Judul Buku} & \textbf{Tahun} & \textbf{Jumlah Halaman} & \textbf{Penerbit} \\
\hline
1 & Struktur Data & 2020 & 90 & Polinema Press \\
\hline
2 & Pemrograman Mobile Menggunakan Flutter & 2021 & 150 & Polinema Press \\
\hline
3 & Laravel REST Api & 2022 & 115 & Polinema Press \\
\hline
\end{longtable}
}

\textbf{G. Perolehan HKI dalam 5 – 10 Tahun Terakhir}

{\footnotesize
\begin{longtable}{|c|p{6cm}|c|c|p{3cm}|}
\hline
\textbf{No.} & \textbf{Judul/Tema HKI} & \textbf{Tahun} & \textbf{Jenis} & \textbf{Nomor P/ID} \\
\hline
1 & Penerapan SQL Similarity Dan Multi-Instance MySQL Database Untuk Proses Penilaian Otomatis Pada Sistem Automated Assessment & 2021 & Program Komputer & EC00202160090 \\
\hline
2 & Sistem Pembelajaran Pronunciation Dan Listening Bahasa Inggris Berbasis Mobile Untuk Sekolah Dasar Kingdem & 2022 & Program Komputer & EC00202297319 \\
\hline
3 & Sistem Informasi Rukun Warga Gondorejo & 2024 & Program Komputer & EC00202482833 \\
\hline
4 & Learning Application In Balanced Assessment (LAIBA) & 2024 & Program Komputer & EC002024204392 \\
\hline
5 & REKACHAIN – Aplikasi Smart Factory Untuk Optimalisasi Lini Produksi & 2024 & Program Komputer & EC002024235519 \\
\hline
\end{longtable}
}

\vspace{0.5cm}

\noindent
Semua data yang saya isikan dan tercantum dalam biodata ini adalah benar dan dapat dipertanggung jawabkan secara hukum. Apabila di kemudian hari ternyata dijumpai ketidaksesuaian dengan kenyataan, saya sanggup menerima sanksi. Demikian biodata ini saya buat dengan sebenarnya untuk memenuhi salah satu persyaratan dalam Penelitian dan Pengabdian Masyarakat.

\vspace{1cm}

\begin{flushright}
Malang, 25 Januari 2026

Pengusul,

\vspace{0.5cm}

\includegraphics[width=3cm]{images/signature-dian.png}

\textbf{Dian Hanifudin Subhi, S.Kom., M.Kom.}

NIP. 198806102019031018
\end{flushright}

\clearpage

\section*{2. Anggota Peneliti 1}

\textbf{A. Identitas Diri}

{\footnotesize
\begin{longtable}{@{}p{0.8cm}p{6cm}p{0.3cm}p{8cm}}
1. & Nama Lengkap (dengan gelar) & : & Dika Rizky Yunianto, S.Kom., M.Kom. \\
2. & Jenis Kelamin & : & L \\
3. & Jabatan Fungsional & : & Asisten Ahli \\
4. & NIK & : & 3507220606920001 \\
5. & NIP & : & 199206062019031017 \\
6. & Tempat dan Tanggal Lahir & : & Malang, 06 Juni 1992 \\
7. & e-mail & : & dika.rizky@polinema.ac.id \\
8. & Nomor Telepon/HP & : & 081944819792 \\
9. & Alamat Kantor & : & Jl. Soekarno Hatta 09, Malang \\
10. & Nomor Telepon/Fax (kantor) & : & (0341) 404424, 404425 \\
11. & Lulusan yg telah dihasilkan & : & 40 orang \\
12. & Mata Kuliah yg diampu & : & Basis Data, Data Warehouse, Data Mining \\
\end{longtable}
}

\textbf{B. Riwayat Pendidikan}

{\footnotesize
\begin{longtable}{|p{3.5cm}|p{5cm}|p{5cm}|p{1.5cm}|}
\hline
& \textbf{S-1} & \textbf{S-2} & \textbf{S-3} \\
\hline
Nama Perguruan Tinggi & Universitas Brawijaya & Institut Teknologi Sepuluh Nopember Surabaya & - \\
\hline
Bidang Ilmu & Teknik Informatika & Teknik Informatika & - \\
\hline
Tahun Masuk-Lulus & 2010-2014 & 2015-2017 & - \\
\hline
Judul Skripsi/Tesis/Disertasi & Segmentasi Citra X-Ray Dental Bitewing Untuk Menentukan Jenis Gigi & Pengembangan Thesaurus Bahasa Arab berbasis Co-Occurence & - \\
\hline
Nama Pembimbing/Promotor & Dr. Candra Dewi, S.Kom., M.Sc & Prof. Dr. Agus Zainal Arifin, S.Kom., M.Kom & - \\
\hline
\end{longtable}
}

\textbf{C. Pengalaman Penelitian Dalam 10 Tahun Terakhir}

{\footnotesize
\begin{longtable}{|c|c|p{6.5cm}|c|c|c|}
\hline
\textbf{No} & \textbf{Tahun} & \textbf{Judul Penelitian} & \textbf{Sebagai} & \multicolumn{2}{c|}{\textbf{Pendanaan}} \\
\cline{5-6}
& & & & \textbf{Sumber} & \textbf{Jumlah (Rp)} \\
\hline
1 & 2023 & ANALISIS SISTEM INFORMASI MONITORING UTILITY PERFORMANCE PADA PT. AIO & Anggota & Politeknik & 8,000,000 \\
\hline
2 & 2023 & APLIKASI REKOMENDASI MAHASISWA BERPRESTASI NON-AKADEMIK DI JURUSAN TEKNOLOGI INFORMASI & Anggota & Politeknik & 8,000,000 \\
\hline
3 & 2022 & PENCARIAN DOKUMEN BERBASIS NORMALISASI LEVENSTHEIN-SIMILARITY PADA APLIKASI JTI MEETING & Anggota & Politeknik & 8,000,000 \\
\hline
4 & 2022 & Inovasi Pembelajaran Hybrid tipe Project Based Learning pada Media Sosial Instagram Marketing & Anggota & Politeknik & 25,000,000 \\
\hline
5 & 2021 & AEGIS APLIKASI CHATBOT CUSTOMER SERVICE JURUSAN TEKNOLOGI INFORMASI & Anggota & Politeknik & 25,000,000 \\
\hline
\end{longtable}
}

\textbf{D. Pengalaman Pengabdian Dalam 10 Tahun Terakhir}

{\footnotesize
\begin{longtable}{|c|c|p{6.5cm}|c|c|c|}
\hline
\textbf{No} & \textbf{Tahun} & \textbf{Judul Penelitian} & \textbf{Sebagai} & \multicolumn{2}{c|}{\textbf{Pendanaan}} \\
\cline{5-6}
& & & & \textbf{Sumber} & \textbf{Jumlah (Rp)} \\
\hline
1 & 2023 & PEMBUATAN SISTEM INFORMASI PENGUMUMAN UNTUK SISWA DAN ORANG TUA DI MIN I KOTA MALANG & Anggota & Politeknik & 7,000,000 \\
\hline
2 & 2022 & SISTEM INFORMASI AKADEMIK BERBASIS WEBSITE DI KB/TK ALFATH BUNULREJO KOTA MALANG & Anggota & Politeknik & 7,000,000 \\
\hline
3 & 2021 & PELATIHAN PEMANFAATAN MICROSOFT OFFICE DAN ADMINISTRATOR OPENSID UNTUK APARATUR DESA TUMPAKOYOT & Anggota & Politeknik & 7,000,000 \\
\hline
\end{longtable}
}

\textbf{E. Pengalaman Publikasi Jurnal Ilmiah}

{\footnotesize
\begin{longtable}{|c|c|p{6cm}|p{3.5cm}|p{2.5cm}|}
\hline
\textbf{No} & \textbf{Tahun} & \textbf{Judul Jurnal Ilmiah} & \textbf{Nama Jurnal Ilmiah} & \textbf{Kategori Publikasi Jurnal Ilmiah} \\
\hline
1 & 2024 & Design And Building Of The Utility Performance Monitoring System At PT. XYZ With Brainstorming Method & International Journal of Health Engineering and Technology & Teknologi Informasi \\
\hline
2 & 2023 & Comparison of Relational Database Modeling Performance based on Number of Normalized Entities & SMARTICS Journal & Teknologi Informasi \\
\hline
3 & 2023 & Natural language processing based chatbot for customer support system on higher education institution & AIP Conference Proceedings & Teknologi Informasi \\
\hline
4 & 2018 & Geo-Sentiment Analysis as a Location-Based Opinion Analysis System on Public Opinion Data about Governor Candidates & International Journal of Engineering \& Technology & Teknologi Informasi \\
\hline
5 & 2018 & Brain signal classification using genetic algorithm for right-left motion pattern & International Journal of Advanced Computer Science and Applications & Teknologi Informasi \\
\hline
6 & 2017 & Pengembangan Sistem Otomatisasi AC dan Lampu Menggunakan Fuzzy dan Raspberry Pi & JURNAL INFOTEL & Teknologi Informasi \\
\hline
7 & 2017 & Perhitungan Kemiripan Term Co-occurence Berdasarkan Cluster Dokumen Untuk Pengembangan Thesaurus Bahasa Arab & JURNAL INFOTEL & Teknologi Informasi \\
\hline
8 & 2015 & Pembobotan Kata BerbasiS Preferensi Dan Hubungan Semantik Pada Dokumen Fiqih Berbahasa Arab & Jurnal Teknologi Informasi dan Ilmu Komputer & Teknologi Informasi \\
\hline
\end{longtable}
}

\textbf{F. Pengalaman HKI Dalam 10 Tahun Terakhir}

{\footnotesize
\begin{longtable}{|c|c|p{8cm}|p{3cm}|}
\hline
\textbf{No} & \textbf{Tahun} & \textbf{Judul} & \textbf{Jenis HKI} \\
\hline
1 & 2020 & Aplikasi Peringkasan Teks Otomatis Portal Berita Online & Hak Cipta \\
\hline
2 & 2023 & Aplikasi Jayanti & Hak Cipta \\
\hline
3 & 2023 & Aplikasi BPSense & Hak Cipta \\
\hline
\end{longtable}
}

\vspace{1cm}

\begin{flushright}
Malang, 22 Januari 2026

\vspace{0.5cm}

\includegraphics[width=3cm]{images/signature-dika.png}

\textbf{Dika Rizky Yunianto, S.Kom., M.Kom.}
\end{flushright}

\clearpage

\section*{3. Anggota Peneliti 2}

\textbf{A. Identitas Diri}

{\footnotesize
\begin{longtable}{@{}p{0.8cm}p{6cm}p{0.3cm}p{8cm}}
1. & Nama Lengkap & : & Ridwan Rismanto, S.ST., M.Kom., Ph.D. \\
2. & Jenis Kelamin & : & Laki-laki \\
3. & Jabatan Fungsional & : & Lektor \\
4. & NIP/NIK/Identitas Lainnya & : & 198603182012121001 \\
5. & NIDN & : & 0018038602 \\
6. & Tempat dan Tanggal Lahir & : & Malang, 18 Maret 1986 \\
7. & E-mail & : & rismanto@polinema.ac.id \\
8. & Nomor Telepon/HP & : & 085646313719 \\
9. & Alamat Kantor & : & Jl. Soekarno-Hatta 9 Malang \\
10. & Nomor Telepon/Faks & : & 0341-404424/0341-404420 \\
11. & Lulusan yang Telah Dihasilkan & : & D3=12; D4=35; S1= - ; S2= - ; S3= - \\
12. & Mata Kuliah yang diampu & : & Praktikum Dasar Pemrograman, Dasar Pemrograman, Struktur Data, Praktikum Struktur Data, Pemrograman Berbasis Objek, Praktikum Pemrograman Berbasis Objek, Metodologi Penelitian, Basis Data, Praktikum Basis Data, Sistem Pendukung Keputusan, Proyek Sistem Informasi, Data Science, Manajemen Korporasi \\
\end{longtable}
}

\textbf{B. Riwayat Pendidikan}

{\footnotesize
\begin{longtable}{|p{3cm}|p{3.5cm}|p{4cm}|p{3.5cm}|}
\hline
& \textbf{S-1} & \textbf{S-2} & \textbf{S-3} \\
\hline
Nama Perguruan Tinggi & Politeknik Elektronika Negeri Surabaya - ITS & Institut Teknologi Sepuluh Nopember & Hiroshima University \\
\hline
Bidang Ilmu & Teknologi Informasi & Teknik Informatika & Informatics and Data Science \\
\hline
Tahun Masuk -Lulus & 2004 – 2008 & 2009-2011 & 2021 - 2025 \\
\hline
Judul Skripsi/Thesis/Disertasi & Perancangan dan Pembuatan Aplikasi Keuangan Berbasis Web & Paralelisasi Transformasi Fourier Pada Arsitektur General Purpose Graphic Processing Unit Untuk Klasifikasi Alat Musik Dengan Instrumen Solo & A Study on Process Evaluation of Concept Map Recomposition \\
\hline
Nama Pembimbing & Isbat Uzzin N, S.Kom., MT. Siyamta, ST., MT. & Dr. Eng. Nanik Suciati, S.Kom., M.Kom. Wahyu Suadi, S.Kom., MM., M.Kom & Prof. Tsukasa Hirashima \\
\hline
\end{longtable}
}

\textbf{C. Pengalaman Penelitian dalam 5 Tahun Terakhir}

{\footnotesize
\begin{longtable}{|c|c|p{10cm}|c|}
\hline
\textbf{No.} & \textbf{Tahun} & \textbf{Judul Penelitian} & \textbf{Sumber} \\
\hline
1 & 2021 & Pengembangan Model Deteksi Lahan Parkir Kosong Menggunakan Metode Deep Learning (tugas belajar) & DIPA \\
\hline
\end{longtable}
}

\textbf{D. Pengalaman Penulisan Artikel Ilmiah dalam Jurnal dalam 5 Tahun Terakhir}

{\footnotesize
\begin{longtable}{|c|p{6cm}|p{3cm}|p{4cm}|}
\hline
\textbf{No.} & \textbf{Judul Artikel Ilmiah} & \textbf{Nama Jurnal} & \textbf{Volume/Nomor/Tahun} \\
\hline
1 & Identifikasi Pneumonia Pada Citra Rontgen Paru Menggunakan Metode Power-Law Trans & Jurnal Informatika Polinema 7 (2), 73-78 & Vol 7/No 2/ 2021 \\
\hline
2 & Sistem UMKM Online Produk Penjualan Kampoeng Sejarah Kota Malang Berbasis Website & Jurnal Pengabdian Polinema Kepada Masyarakat 8 (2), 102-104 & Vol 8/No 2/2021 \\
\hline
3 & Analysis of Comparison of Fuzzy Knn, C4. 5 Algorithm, and Naïve Bayes Classification Method for Diabetes Mellitus Diagnosis & International Journal of Computer Applications Technology and Research & Vol 7/No 8/2023 \\
\hline
4 & Experimental Evaluation of OOP Concepts Learning with a Kit-Build Concept Map: A Case Study & International Journal of Information and Education Technology 13 (9), 1358-1364 & Vol 13/No 9/2023 \\
\hline
5 & Evaluating the kit-build concept mapping process using sub-map scoring & Research and Practice in Technology Enhanced Learning & Vol 19/No 21/2024 \\
\hline
6 & The evaluation of concept map recomposition order and its effect on learning outcomes & Research and Practice in Technology Enhanced Learning & Vol 21/2026 \\
\hline
\end{longtable}
}

\textbf{E. Pemakalah Seminar Ilmiah (Oral Presentation) dalam 5 Tahun Terakhir}

{\footnotesize
\begin{longtable}{|c|p{5.5cm}|p{5cm}|p{3.5cm}|}
\hline
\textbf{No.} & \textbf{Nama Pertemuan Ilmiah/Seminar} & \textbf{Judul Artikel} & \textbf{Waktu dan Tempat} \\
\hline
1 & A Preliminary Study: Toulmin Arguments in English Reading Comprehension for English as Foreign Language Students & 2021 International Conference on Electrical and Information Technology & September 2021, Malang, Indonesia \\
\hline
2 & Rekomendasi e-Book pada Digital Repository Rumah Sakit Umum Universitas Kristen Indonesia dengan menggunakan Metode VIKOR & Seminar Informatika Aplikatif Polinema & November 2021, Malang, Indonesia \\
\hline
3 & Process Evaluation for Concept Map Building and Its Experimental Evaluation & International Conference on Computers in Education & Desember 2023, Shimane, Jepang \\
\hline
4 & Trial-and-Error Modifications as Cognitive Compensation: Analyzing the Concept Mapping Process in Reconstruction Tasks & International Conference on Computers in Education & Desember 2025, Chennai, India \\
\hline
\end{longtable}
}

\textbf{F. Perolehan HKI dalam 5 – 10 Tahun Terakhir}

{\footnotesize
\begin{longtable}{|c|p{8cm}|c|c|c|}
\hline
\textbf{No.} & \textbf{Judul/Tema HKI} & \textbf{Tahun} & \textbf{Jenis} & \textbf{Nomor P/ID} \\
\hline
1 & Program Rekomendasi Artikel Terkait Pada Berita Online Menggunakan Teknik Text Mining (tugas belajar) & 2019 & Program Komputer & EC00201977365 \\
\hline
\end{longtable}
}

\vspace{0.5cm}

\noindent
Semua data yang saya isikan dan tercantum dalam biodata ini adalah benar dan dapat dipertanggung jawabkan secara hukum. Apabila di kemudian hari ternyata dijumpai ketidaksesuaian dengan kenyataan, saya sanggup menerima sanksi. Demikian biodata ini saya buat dengan sebenarnya untuk memenuhi salah satu persyaratan dalam Penelitian dan Pengabdian Masyarakat.

\vspace{1cm}

\begin{flushright}
Malang, 25 Januari 2026

\vspace{0.5cm}

\includegraphics[width=3cm]{images/signature-ridwan.png}

\textbf{Ridwan Rismanto, S.ST., M.Kom., Ph.D.}

NIP. 198603182012121001
\end{flushright}

\clearpage

\section*{4. Anggota Peneliti 3}

\textbf{A. Identitas Diri}

{\footnotesize
\begin{longtable}{@{}p{0.8cm}p{6cm}p{0.3cm}p{8cm}}
1. & Nama Lengkap (dengan gelar) & : & Usman Nurhasan, S.Kom., M.T \\
2. & Jenis Kelamin & : & Laki-laki \\
3. & Jabatan Fungsional & : & Lektor \\
4. & NIP/NIK/No. identitas lainnya & : & 198609232015041001 \\
5. & NIDN & : & 0023098604 \\
6. & Tempat dan Tanggal Lahir & : & Malang, 23 September 1986 \\
7. & E-mail & : & usmannurhasan@polinema.ac.id \\
8. & Nomor Telepon/HP & : & 082143246664 \\
9. & Alamat Kantor & : & Jl. Soekarno Hatta no. 9 Malang 65141 \\
10. & Nomor Telepon/Faks & : & 0341-404424 \\
11. & Mata Kuliah yg Diampu & : & Multimedia Terapan, Pemrograman Jaringan, Aplikasi Enterprise, Jaringan Komputer, Sistem Operasi \\
\end{longtable}
}

\textbf{B. Riwayat Pendidikan}

{\footnotesize
\begin{longtable}{|p{3.5cm}|p{5.5cm}|p{6cm}|}
\hline
& \textbf{S-1} & \textbf{S-2} \\
\hline
Nama Perguruan Tinggi & Universitas Islam Negeri Maulana Malik Ibrahim Malang & Universitas Brawijaya Malang \\
\hline
Bidang Ilmu & Teknik Informatika & Teknik Elektro Minat Sistem Komunikasi dan Informatika \\
\hline
Tahun Masuk- Lulus & 2005-2010 & 2011-2014 \\
\hline
Judul Skripsi/ Tesis/Disertasi & Rancang Bangun Aplikasi Pengontrol Peralatan Listrik Menggunakan NM7010 Adan DT-AVR Low Cost Series & Sistem Cerdas Pencarian Lokasi Parkir Terbaik dengan Algoritma Fuzzy-Best Fit \\
\hline
Nama Pembimbing/ Promotor & 1. Totok Chamidy, S.Kom., M.Kom 2. Dr. Ahmad Barizy, M.Ag & 1. Hadi Suyono, ST., MT., Ph.D 2. Dr. Aziz Muslim, MT \\
\hline
\end{longtable}
}

\textbf{C. Pengalaman Penelitian Dalam 5 Tahun Terakhir}

{\footnotesize
\begin{longtable}{|c|c|p{6cm}|c|c|}
\hline
\textbf{No.} & \textbf{Tahun} & \textbf{Judul Penelitian} & \textbf{Sumber*} & \textbf{Jml (Juta Rp)} \\
\hline
1 & 2016 & Sistem Cerdas Tes Psikologi Papi Kostick Untuk Seleksi Masuk Mahasiswa Baru Jalur D-III Kerjasama Polinema Dengan PLN Dan PT GMF Aeroasia & DIPA & 4 \\
\hline
2 & 2016 & Pengembangan Sistem Analisa Sentimen Pada Kata Kunci Tertentu Di Dalam Sebuah Korpus Teks & DIPA & 4 \\
\hline
3 & 2017 & Implementasi Fuzzy Inference System Pada Sistem Cerdas Tes Psikologi Kraepelin & DIPA & 4 \\
\hline
4 & 2017 & Sistem Cerdas Penentu Rekomendasi Pemilihan Jalur Angkot Kota Malang & DIPA & 4 \\
\hline
5 & 2018 & Efektifitas Mobile Art Application (MATA) Sebagai Media Penanganan Kasus Agresi Pada Anak & DIPA & 4 \\
\hline
6 & 2018 & Implementasi Teksturing Secara Real Time Pada Augmented Reality (AR) & DIPA & 4 \\
\hline
7 & 2018 & Implementasi Smart Mirror Tracer Untuk Mendeteksi Tingkat Dominan dan Emosi & DIPA & 25 \\
\hline
8 & 2019 & Penerapan Teknologi Hybrid (Augmented Reality Dan Geographic Information System) Untuk Tourist Guide Application & DIPA & 8 \\
\hline
9 & 2019 & Perancangan Game Edukasi Pengenalan Huruf Hijaiyah Untuk Anak Berbasis Augmented Reality & DIPA & 8 \\
\hline
10 & 2019 & Aplikasi Pembelajaran Bahasa Inggris Berbasis Augmented Reality & DIPA & 8 \\
\hline
11 & 2019 & Pemanfaatan Konten Virtual Reality Untuk Simulasi Matakuliah Praktikum & DIPA & 50 \\
\hline
12 & 2020 & Implementasi E-Reporting Infrastructure (Eri) Sebagai Media Monitoring Kerusakan Infrastruktur Dalam Kampus & DIPA & 8 \\
\hline
13 & 2020 & Smart Ecosystem Untuk Lahan Hidroponik Pada Kelompok Petani Hidroponik Binaan CSR PT. Otsuka Indonesia Sebagai Peningkatan Kualitas Dan Kuantitas Hasil Panen & DIPA & 70 \\
\hline
14 & 2021 & Implementasi Sistem Cerdas Eksperimen Gambar Untuk Analisis Kondisi Adversity Response Profile (ARP) & DIPA & 8 \\
\hline
\end{longtable}
}

\textbf{D. Pengalaman Pengabdian Kepada Masyarakat dalam 5 Tahun Terakhir}

{\footnotesize
\begin{longtable}{|c|c|p{6cm}|c|c|}
\hline
\textbf{No.} & \textbf{Tahun} & \textbf{Judul Pengabdian Kepada Masyarakat} & \textbf{Sumber*} & \textbf{Jml (Juta Rp)} \\
\hline
1. & 2016 & Penerapan Sistem Pengawasan IP Camera di Yayasan Pendidikan Islam & DIPA & 3,6 \\
\hline
2 & 2017 & Pelatihan Pengolahan Nilai Raport kepada Guru di SDN 1 Sukun Kota Malang & DIPA & 3,6 \\
\hline
3 & 2018 & Packaging dan Promosi Produk Bahan Olahan Makanan Untuk Meningkatkan Daya Saing (Studi Kasus UKM "Makanan Olahan Menuk") & DIPA & 3,6 \\
\hline
4 & 2019 & Pelatihan Pemanfaatan Aplikasi Smart.04 & DIPA & 6,8 \\
\hline
5 & 2019 & Pelatihan Keterampilan Dasar Komputer Dan Pengolahan Nilai Raport Bagi Guru Smp Islam Al-Amin Kota Malang & DIPA & 6,8 \\
\hline
6 & 2020 & Pelatihan Pemanfaatan Marketplace Tani-Ku Sebagai Media Promosi Dan Transaksi Bagi Penggiat Hidroponik Di Kota Batu & DIPA & 6,8 \\
\hline
7 & 2020 & Pelatihan Keterampilan Dasar Komputer Dan Tata Kelola Organisasi Pada Komunitas Pemuda Penggiat Sejarah Di Kelurahan Sumbersari Kota Malang & DIPA & 6,8 \\
\hline
8 & 2020 & Sosialisasi Pentingnya Literasi Digital Untuk Meningkatkan Komunikasi Digital Bagi Komunitas Reenactor Ngalam & DIPA & 6,8 \\
\hline
9 & 2021 & Pelatihan Pemanfaatan Media Literasi Mandiri Di Lingkungan Madrasah Aliyah Negeri 1 Malang Sebagai Usaha Mewujudkan Smart School & DIPA & 7 \\
\hline
\end{longtable}
}

\textbf{E. Publikasi Artikel Ilmiah Dalam Jurnal dalam 5 Tahun Terakhir}

{\footnotesize
\begin{longtable}{|c|p{5.5cm}|p{5cm}|p{3cm}|}
\hline
\textbf{No.} & \textbf{Judul Artikel Ilmiah} & \textbf{Nama Jurnal} & \textbf{Volume/ Nomor/Tahun} \\
\hline
1 & Sistem Cerdas Pencarian Lokasi Parkir Terbaik dengan Algoritma Fuzzy-Best Fit & Electrical Power, Electronics, Communications, Controls, and Informatics (EECCIS) & Volume 8 / No 1 / Tahun 2014 \\
\hline
2 & Sistem Cerdas Papi Kostick & JURNAL TEKNO LOGI INFORMASI: Teori, Konsep, dan Implementasi ISSN : 2086 - 2989 & Volume 12/ No 1 /Tahun 2017 \\
\hline
3 & Sistem Cerdas Tes Kepribadian Kraepelin & JURNAL TEKNO LOGI INFORMASI: Teori, Konsep, dan Implementasi ISSN : 2086 - 2989 & Volume 8 / No. 2/ 2017 \\
\hline
4 & Sistem Cerdas Penentu Rekomendasi Pemilihan Jalur Angkot Kota Malang & Jurnal ANTIVIRUS :Jurnal Ilmiah Teknik Informatika ISSN 1978-5232 & Volume 11/ Tahun 2017 \\
\hline
5 & Implementation IoT in System Monitoring Hydroponic Plant Water Circulation and Control & International Journal of Engineering and Technology(UAE) ISSN: 2227-524X & Volume 7/2018 \\
\hline
6 & Realtime Mobile Queue System Multi Vendor & Jurnal Informatika Polinema ISSN: 2614-6371 & Volume 5/Tahun 2019 \\
\hline
7 & Pemanfaatan Website Untuk Manajemen Informasi Penelitian Peternakan Kota Batu & Jurnal Antivirus: Jurnal Ilmiah Teknik Informatika ISSN 1978-5232 & Volume 13/Tahun 2019 \\
\hline
8 & Perancangan Sistem Informasi Pengolahan Data Belanja Langsung Berbasis Website & Generation Journal e-ISSN number: 2549-2233 p-ISSN: 2580-4952 & Volume 3 / Tahun 2019 \\
\hline
9 & Simulasi Virtual Reality Pendeteksi Tingkat Phobia Seseorang Sebagai Alat Bantu Terapi Acrophobia (Phobia Ketinggian) Berbasis Android & Jurnal Informatika Polinema ISSN: 2614-6371 & Volume 5/Tahun 2019 \\
\hline
10 & Sistem Informasi Geografis Pencarian Jalur Terdekat Desa Wisata Studi Kasus Kota Wisata Batu & Jurnal Informatika Polinema ISSN: 2614-6371 & Volume 5/Tahun 2019 \\
\hline
11 & Mobile Art Application (MATA) as a media for handling cases of aggression in children & Journal of Physics: Conference Series & 1402 (6), 066049, 2019 \\
\hline
12 & Study of image reconstructions of Indonesian traffic sign using Criminisi and Hung method & Journal of Physics: Conference Series & 1402 (6), 066032 \\
\hline
13 & "Hijaiyah" interactive learning for pre-school students & Journal of Physics: Conference Series & 1402 (6), 066050 \\
\hline
14 & Smart tourist guide application for the introduction of Malang city tourism potential using hybrid technology & Journal of Physics: Conference Series & 1402 (6), 066051 \\
\hline
15 & Utilization of virtual reality content for laboratory practicum learning & IOP Conference Series: Materials Science and Engineering & 732 (1), 012087/ 2020 \\
\hline
16 & Pengembangan Aplikasi Deteksi Dini Pada Anak Berkebutuhan Khusus Di Sekolah Luar Biasa & Petir: Jurnal Pengkajian dan Penerapan Teknik Informatika & 14 (1), 17-27/ 2020 \\
\hline
17 & Implementasi Metode MVC Untuk Pembangunan Sistem Informasi Pelatihan Kerja: Studi Kasus UPT Pelatihan Kerja Singosari & Jurnal Eksplora Informatika & 10 (1), 20-31 1/ 2020 \\
\hline
18 & Pembelajaran Humane Education Di Kurikulum 2013 Menggunakan Media Animasi & Urbangreen Conference Proceeding Library & 1, 176-182/ 2020 \\
\hline
19 & Analisis Perilaku Non Playable Character (Npc) Pada Game Menggunakan Fuzzy Sugeno & Techno. Com & 19 (3), 308-320/ 2020 \\
\hline
20 & Penerapan Aplikasi Up-Event Pada Tata Kelola Kegiatan Multi Vendor & Generation Journal & 4 (2), 84-97 \\
\hline
21 & Pemodelan Konten Augmented Reality Dalam Konsep Interactive Realtime & Jurnal Informatika & 20 (1), 30-40/2020 \\
\hline
22 & Media Pembelajaran Interaktif Terhadap Gwk Culture Park Menggunakan Teknologi Augmented Reality & Jurnal Informatika Polinema & 6 (3), 1-8/2020 \\
\hline
23 & Implementation of Hybrid Methods in the Application of Experimental Psychology for Analysis of Mental Endurance Conditions & Khazanah Informatika: Jurnal Ilmu Komputer dan Informatika 7 (1) & Vol 7/ No 1/2021 \\
\hline
24 & Mobile Media Literature Application For Identification Of Student's Reading Interest & Jurnal Riset Informatika (JRI) & Vol 3 No 4 (2021): Period of September 2021 \\
\hline
25 & Penerapan Vector Space Model Dalam Klasifikasi Penilaian Thematic Appeception Test & J-SAKTI (Jurnal Sains Komputer \& Informatika) & Vol 5, No 2 / 2021 \\
\hline
26 & Implementasi Metode Fuzzy Time Series Markov Chain untuk Prediksi Harga Telur Puyuh & Informatika Mulawarman : Jurnal Ilmiah Ilmu Komputer & Vol 16, No 2/2021 \\
\hline
\end{longtable}
}

\textbf{F. Pemakalah Seminar Ilmiah (Oral Presentation) dalam 5 Tahun Terakhir}

{\footnotesize
\begin{longtable}{|c|p{5cm}|p{5.5cm}|p{3.5cm}|}
\hline
\textbf{No.} & \textbf{Nama Temu ilmiah/Seminar} & \textbf{Judul Artikel Ilmiah} & \textbf{Waktu dan Tempat} \\
\hline
1 & MATEC Web of Conferences & The evaluation on a mobile augmented reality application as therapy media & 2018, Bandung, Jawa Barat, Indonesia \\
\hline
2 & Proceedings of the 2018 International Conference on Energy and Mining Law (ICEML 2018) & Analysis of Sand Mining Areas in Lumajang Using WEBGIS & 2018, Bandung, Jawa Barat, Indonesia \\
\hline
3 & The 4th Annual Applied Science and Engineering Conference & Mobile Art Application (MATA) as a media for handling cases of aggression in children & 2019, Bali, Indonesia \\
\hline
4 & The 4th Annual Applied Science and Engineering Conference & Study of image reconstructions of Indonesian traffic sign using Criminisi and Hung method & 2019, Bali, Indonesia \\
\hline
5 & The 4th Annual Applied Science and Engineering Conference & "Hijaiyah" interactive learning for pre-school students & 2019, Bali, Indonesia \\
\hline
6 & The 4th Annual Applied Science and Engineering Conference & Smart tourist guide application for the introduction of Malang city tourism potential using hybrid technology & 2019, Bali, Indonesia \\
\hline
7 & Seminar Nasional Terapan Riset Inovatif (Sentrinov) 2019 & Penerapan Socket Programming Pada Game Multi Client "Di Kejar Musuh" & 2019, Bangka Belitung, Indonesia \\
\hline
8 & Seminar Nasional Terapan Riset Inovatif (Sentrinov) 2019 & Pemanfaatan Konten Virtual Reality Sebagai Media Terapi Penderita Ophidiophobia & 2019, Bangka Belitung, Indonesia \\
\hline
9 & The 2nd International Conference on Applied Science and Technology (iCAST) 2019, & Implementation of Augmented Reality For Autism Children Therapy Spectrum Disorder & 2019, Bali, Indonesia \\
\hline
10 & The 2nd Annual Technology, Applied Science and Engineering Conference (ATASEC 2020) & Smart Ecosystem For Hydroponic Land In The Hydroponic Farmers Group Guided By CSR PT. Otsuka Indonesia As An Improved Quality And Quantity Of Harvest Results & 2020, Indonesia \\
\hline
11 & Seminar Nasional Psikologi UM (2021) & Problematika Pembelajaran Jarak Jauh (PJJ) Mata Pelajaran Informatika pada Masa Pandemi di MAN 1 Kota Malang & 2021, Malang \\
\hline
12 & UrbanGreen Conference Proceeding Library & Pembelajaran Humane Education di Kurikulum 2013 Menggunakan Media Animasi & 2021, Malang \\
\hline
13 & International Conference on Industrial Revolution for Polytechnic Education & Application of WebSocket on Integration of E-Reporting Infrastructure (ERI) Applications with IoT & 2021, Malang \\
\hline
\end{longtable}
}

\textbf{G. Karya Buku dalam 5 Tahun Terakhir}

{\footnotesize
\begin{longtable}{|c|p{7cm}|c|c|p{2.5cm}|}
\hline
\textbf{No.} & \textbf{Judul Buku} & \textbf{Tahun} & \textbf{Jumlah Halaman} & \textbf{Penerbit} \\
\hline
1 & Pemrograman Jaringan Menggunakan Java Jilid 1 & 2018 & 107 & Polinema Press \\
\hline
2 & Pemrograman Jaringan Menggunakan Java Jilid 2 & 2018 & 260 & Polinema Press \\
\hline
3 & Isu Profesional dan Sosial di Bidang IT & 2019 & 132 & Polinema Press \\
\hline
4 & Pengantar Sistem Operasi dan Remastering Linux & 2020 & 174 & Polinema Press \\
\hline
5 & Pengantar Rekayasa Perangkat Lunak & 2020 & 190 & Polinema Press \\
\hline
\end{longtable}
}

\textbf{H. Perolehan HKI dalam 10 Tahun Terakhir}

{\footnotesize
\begin{longtable}{|c|p{6cm}|c|c|c|}
\hline
\textbf{No.} & \textbf{Judul/Tema HKI} & \textbf{Tahun} & \textbf{Jenis} & \textbf{Nomor P/ID} \\
\hline
1 & ARDoa Augmented Reality Doa & 2018 & Hak Cipta & EC00201858687 \\
\hline
2 & Media Pembelajaran Berbasis Animasi "Animal Education" & 2019 & Hak Cipta & EC00201983866 \\
\hline
3 & E-Reporting Infrastruktur (ERI) Kampus & 2020 & Hak Cipta & EC00202044987 \\
\hline
4 & Aplikasi Marketplace Tani-Ku & 2020 & Hak Cipta & EC00202044986 \\
\hline
5 & Art Drawing Sebagai Media Komunikasi Nonverbal & 2020 & Hak Cipta & EC00202044046 \\
\hline
6 & Aplikasi E-Payment Bank Sampah & 2020 & Hak Cipta & EC00202044045 \\
\hline
7 & ARSurya : Aplikasi Media Pembelajaran Tat Surya Berbasis Augmented Reality & 2020 & Hak Cipta & EC00202044044 \\
\hline
\end{longtable}
}

\vspace{0.5cm}

\noindent
Semua data yang saya isikan dan tercantum dalam biodata ini adalah benar dan dapat dipertanggungjawabkan secara hukum. Apabila di kemudian hari ternyata dijumpai ketidak-sesuaian dengan kenyataan, saya sanggup menerima risikonya.

\noindent
Demikian biodata ini saya buat dengan sebenarnya untuk memenuhi salah satu persyaratan dalam pengajuan penelitian dan pengabdian.

\vspace{1cm}

\begin{flushright}
Malang, 25 Januari 2026

\vspace{0.5cm}

\includegraphics[width=3cm]{images/signature-usman.png}

\textbf{Usman Nurhasan, S.Kom., M.T}

NIP. 198609862015041001
\end{flushright}
